%% kapa: exam=1 round=30 year=2019 period=1 form=A subject=civil_law subject_ko=민법 seq=33 number=33
%% source: Q-Net 감정평가사 기출문제 (cst00309) artlSeq=5205205 file=extracted/1008/5205205/2019년 제30회 감정평가사 1차시험 문제지/2019년 제30회 감정평가사 1차시험 문제지/2019년도 제30회 감정평가사 1차 1교시 A형.hwp
%% conversion: kordoc 4.24.51 -> tools/gen1.py (auto); review=proofread-sonnet; flags=none
\begin{question}{33}
2017년 8월경 甲은 乙소유의 X부동산을 매매대금을 일시에 지급하고 매수하면서 애인인 丙과의 명의신탁약정에 기초하여 乙로부터 丙으로 X부동산에 관한 소유권이전등기를 마쳤다. 이에 관한 설명으로 옳지 않은 것은? (다툼이 있으면 판례에 따름)
\begin{choices}
  \item 甲과 丙 사이의 명의신탁약정 및 그에 따른 丙 명의의 등기는 무효이다.
  \item 甲과 丙이 이후 혼인을 하게 된다면, 조세포탈 등이나 법령상의 제한을 회피할 목적이 없는 한, 위 등기는 甲과 丙이 혼인한 때로부터 유효하게 된다.
  \item 丙이 X부동산을 임의로 처분하였다 하더라도 특별한 사정이 없는 한, 乙이 丙의 처분행위로 인하여 손해를 입었다고 할 수는 없다.
  \item 甲은 乙에 대한 소유권이전등기청구권을 보전하기 위하여 乙을 대위하여 丙 명의의 등기말소를 청구할 수 있다.
  \item 丙으로부터 X부동산을 매수한 丁이 丙의 甲에 대한 배임행위에 적극 가담하였더라도, 丙과 丁 사이의 매매계약은 반사회적인 법률행위에 해당하지는 않는다.
\end{choices}
\end{question}
\answer{5}
